\documentclass[11pt,a4paper]{article}
\usepackage[utf8]{inputenc}
\usepackage[T1]{fontenc}
\usepackage[spanish]{babel}
\usepackage{amsmath,amssymb}
\usepackage{booktabs}
\usepackage{geometry}
\usepackage{enumitem}
\usepackage{hyperref}
\usepackage{xcolor}

\geometry{margin=2.5cm}
\hypersetup{colorlinks=true, linkcolor=blue!60!black, urlcolor=blue!60!black}

\newcommand{\metrica}[1]{\textbf{\textsf{#1}}}
\newcommand{\pp}{\,\text{pp}}

\title{\vspace{-1.5cm}\textbf{Protocolo de Evaluaci\'on: Structured Pool}\\[4pt]
\large M\'etricas del sistema cross-modal Audio\,$\leftrightarrow$\,MIDI\\[2pt]
\normalsize Proyecto Phideus --- BIAS\_CONTROL}
\author{}
\date{}

\begin{document}
\maketitle
\thispagestyle{empty}

% ============================================================
\section{El setup: qu\'e se eval\'ua}

El modelo genera embeddings de 256 dimensiones.
Dado un segmento de audio produce un vector $\mathbf{e}_a \in \mathbb{R}^{256}$;
dado un segmento MIDI produce $\mathbf{e}_m \in \mathbb{R}^{256}$.
Si el modelo aprendi\'o correctamente, los vectores de audio y MIDI
\emph{del mismo segmento} deber\'ian quedar cercanos en ese espacio.

Para testear eso se arma un \textbf{pool estructurado} de 256 candidatos
por cada query, dise\~nados para ser dif\'iciles:

\begin{table}[h]
\centering
\begin{tabular}{@{}lrl@{}}
\toprule
\textbf{Tipo} & \textbf{Cantidad} & \textbf{Descripci\'on} \\
\midrule
Positivo        & 1   & El segmento correcto (misma pieza, mismo momento) \\
Hard negatives  & 64  & Misma pieza, distinto momento --- lo m\'as dif\'icil \\
Semi-hard       & 32  & Mismo compositor, otra pieza \\
Random          & 159 & Cualquier otra cosa \\
\midrule
\textbf{Total}  & \textbf{256} & \\
\bottomrule
\end{tabular}
\caption{Composici\'on del pool estructurado por query.}
\label{tab:pool}
\end{table}

Se ejecutan $Q = 500$ queries con \texttt{seed=42}
(siempre las mismas, para comparabilidad entre runs).

% ============================================================
\section{\metrica{A2M R@10} --- Audio-to-MIDI Recall at 10}

\textbf{Pregunta que responde:}
\emph{``Si te doy un audio, \textquestiondown encontr\'as su MIDI correcto
dentro de tus 10 mejores candidatos?''}

\subsection*{C\'omo se calcula}

\begin{enumerate}[nosep]
  \item Tomar el embedding del audio query: $\mathbf{e}_a^{(q)}$.
  \item Calcular la similitud coseno contra los 256 embeddings MIDI del pool:
    \[
      s_i = \frac{\mathbf{e}_a^{(q)} \cdot \mathbf{e}_m^{(i)}}
                 {\|\mathbf{e}_a^{(q)}\|\;\|\mathbf{e}_m^{(i)}\|},
      \qquad i = 1, \dots, 256.
    \]
  \item Ordenar de mayor a menor similitud.
  \item Definir el indicador por query:
    \[
      h(q) =
      \begin{cases}
        1 & \text{si el MIDI correcto est\'a en los top-10,} \\
        0 & \text{en caso contrario.}
      \end{cases}
    \]
  \item Promediar sobre las $Q$ queries:
    \[
      \boxed{\;\text{A2M R@10} = \frac{1}{Q}\sum_{q=1}^{Q} h(q)\;}
    \]
\end{enumerate}

\subsection*{Interpretaci\'on}

Si $\text{A2M} = 60.2\%$, significa que en 301 de 500 queries el modelo
encontr\'o el MIDI correcto entre sus 10 mejores picks, enfrentando
246 candidatos distractores (incluyendo 64 de la misma pieza).

% ============================================================
\section{\metrica{M2A R@10} --- MIDI-to-Audio Recall at 10}

Exactamente lo mismo pero en direcci\'on inversa: se entrega un embedding MIDI
como query y se busca su audio correcto en el pool de 256 embeddings de audio.

\[
  \text{M2A R@10} = \frac{1}{Q}\sum_{q=1}^{Q} h'(q),
  \qquad
  h'(q) =
  \begin{cases}
    1 & \text{si el audio correcto est\'a en top-10,} \\
    0 & \text{en caso contrario.}
  \end{cases}
\]

\subsection*{Por qu\'e importa que sean dos m\'etricas separadas}

Si el modelo aprendi\'o una direcci\'on pero no la otra
(e.g.\ ``dado audio encuentro MIDI'' pero ``dado MIDI no encuentro audio''),
algo est\'a desbalanceado.
Esto fue exactamente lo que ocurri\'o en Gate~4/4.1:
M2A sub\'ia mientras A2M se degradaba,
porque el audio encoder estaba congelado
y solo se mov\'ia el MIDI encoder.

% ============================================================
\section{\metrica{S} $= \min(\text{A2M},\;\text{M2A})$ --- Score balanceado}

\textbf{Pregunta que responde:}
\emph{``\textquestiondown Cu\'al es el peor de los dos sentidos?''}

\[
  \boxed{\;S = \min\!\big(\text{A2M R@10},\;\text{M2A R@10}\big)\;}
\]

Es la \textbf{m\'etrica primaria de decisi\'on}.
Se eligi\'o $\min$ en vez de promedio para \emph{castigar asimetr\'ias}.

\begin{itemize}[nosep]
  \item Si $\text{A2M} = 80\%$ y $\text{M2A} = 30\%$:
    \begin{itemize}[nosep]
      \item Promedio $= 55\%$ \quad (parece bien).
      \item $S = 30\%$ \quad (refleja que una direcci\'on est\'a rota).
    \end{itemize}
\end{itemize}

$S$ es la m\'etrica que define \textsc{go}/\textsc{no-go} para cada run.

% ============================================================
\section{\metrica{hard\_neg} --- Hard Negative Accuracy}

\textbf{Pregunta que responde:}
\emph{``\textquestiondown Puede el modelo distinguir el segmento correcto
de otros segmentos \textbf{de la misma pieza}?''}

Este es el test m\'as exigente.
Los 64 hard negatives son segmentos de la misma pieza musical
pero en momentos diferentes.
El audio suena parecido
(mismo piano, mismo compositor, misma sala, mismo estilo).
La \'unica diferencia son las notas exactas que suenan en ese instante.

\subsection*{C\'omo se calcula}

\begin{enumerate}[nosep]
  \item Para cada query, considerar solo los hard negatives $+$ el positivo
        (65 candidatos).
  \item Definir:
    \[
      g(q) =
      \begin{cases}
        1 & \text{si el positivo tiene mayor similitud que los 64 hard negatives,} \\
        0 & \text{en caso contrario.}
      \end{cases}
    \]
  \item Promediar:
    \[
      \boxed{\;\text{hard\_neg} = \frac{1}{Q}\sum_{q=1}^{Q} g(q)\;}
    \]
\end{enumerate}

\subsection*{Interpretaci\'on}

Si $\text{hard\_neg} = 91\%$, el modelo distingue
\emph{``este momento exacto de esta Sonata de Beethoven''}
de \emph{``otro momento de esa misma Sonata''}
en 455 de 500 queries.

\subsection*{Por qu\'e importa}

Un modelo podr\'ia tener A2M alto solo por aprender
\emph{firma de pieza}
(reconoce que es la Sonata~14 de Beethoven,
pero no sabe \emph{qu\'e parte}).
hard\_neg detecta eso:
si solo aprendi\'o firma,
$\text{hard\_neg} \approx \frac{1}{65} \approx 1.5\%$.

% ============================================================
\section{\metrica{|gap|} $= |\text{A2M} - \text{M2A}|$ --- Asimetr\'ia}

\textbf{Pregunta que responde:}
\emph{``\textquestiondown El modelo funciona igual de bien en ambas direcciones?''}

\[
  \boxed{\;|\text{gap}| = \big|\text{A2M R@10} - \text{M2A R@10}\big|\;}
\]

\begin{itemize}[nosep]
  \item $|\text{gap}|$ bajo (${\sim}0$--$2\pp$):
        simetr\'ia saludable, ambos encoders contribuyen.
  \item $|\text{gap}|$ alto (${\sim}5$--$7\pp$):
        un sentido domina; probable que un encoder est\'e subentrenado.
\end{itemize}

En Run~D original, $|\text{gap}|$ baj\'o a $0.8\pp$ en epoch~5
--- la mejor simetr\'ia de todos los runs.
En D-02 oscila entre $0.4\pp$ y $3.4\pp$ en la meseta,
con mediana ${\sim}1.6\pp$.

% ============================================================
\section{Ejemplo concreto: D-02 epoch~18}

\[
  S = 59.6\%,\quad
  \text{A2M} = 60.8\%,\quad
  \text{M2A} = 59.6\%,\quad
  \text{hard\_neg} = 91.0\%,\quad
  |\text{gap}| = 1.2\pp
\]

\textbf{Traducci\'on:}
De 500 queries audio$\to$MIDI, el modelo ubica el MIDI correcto en top-10
el $60.8\%$ de las veces.
Al rev\'es (MIDI$\to$audio), $59.6\%$.
El peor caso es $59.6\%$ (M2A limita, por eso $S = 59.6\%$).
Contra segmentos de la misma pieza, acierta el $91\%$ de las veces.
La diferencia entre ambas direcciones es solo $1.2$ puntos porcentuales:
el modelo est\'a balanceado.

% ============================================================
\section{Comparativa para dimensionar}

\begin{table}[h]
\centering
\begin{tabular}{@{}lccl@{}}
\toprule
\textbf{Referencia} & $\boldsymbol{S}$ & \textbf{hard\_neg} & \textbf{Significado} \\
\midrule
Random              & $3.9\%$   & $1.5\%$   & Azar puro \\
Gate~2 (baseline)   & $34.4\%$  & $80.4\%$  & Primer modelo funcional \\
Run~D (5 epochs)    & $51.0\%$  & $89.2\%$  & Foundation provisional \\
\textbf{D-02 (e18)} & $\boldsymbol{59.6\%}$ & $\boldsymbol{91.0\%}$ & \textbf{Mejor hasta ahora} \\
Techo te\'orico     & $100\%$   & $100\%$   & Matching perfecto \\
\bottomrule
\end{tabular}
\caption{Comparativa de m\'etricas a trav\'es de las etapas del proyecto.}
\label{tab:comparativa}
\end{table}

\end{document}
